\section{模型建立}

\subsection{右删失物理信息排放数字孪生}

\subsubsection{Deutsch--Anderson 穿透骨架}

经典 Deutsch--Anderson 模型以指数形式描述电除尘器穿透率\cite{nichols1972epa,white1963esp}。对四电场设备，定义无振打基线浓度
\begin{equation}
C_{\mathrm{base},t}=1000C_{\mathrm{in},t}\exp(-D_t),
\label{eq:basec}
\end{equation}
其中 $1000$ 用于将 $\gNm$ 转为 $\mgNm$，等效去除指数写为
\begin{equation}
D_t=D_{0,t}\frac{Q_0}{Q_t}\,g_\Theta(\Theta_t)\,
\underbrace{\sum_{i=1}^{4}w_{i,t}\left(\frac{U_{i,t}}{U_{i,0}}\right)^2}_{S_U(\bm U_t)}
\,g_M(\bm T_t,L_t).
\label{eq:dindex}
\end{equation}
式中 $\sum_iw_{i,t}=1$，$Q_0,U_{i,0}$ 为数据中位参考点。$Q_0/Q_t$ 对应停留时间；$U_i^2$ 对应迁移驱动力的低阶近似。附件没有分场出口浓度、二次电流和极板面积，无法独立回归四场效率。为避免手工指定，在四场未观测几何与迁移系数等效的最小信息假设下，令参考点各场贡献与 $U_{i,0}^2$ 成正比：
\begin{equation}
 w_{i,0}=\frac{U_{i,0}^2}{\sum_{j=1}^{4}U_{j,0}^2},\qquad
\bm w_0=(0.2978,0.2988,0.2017,0.2017),
\label{eq:weights}
\end{equation}
其中 $\bm U_0=(58.8,58.9,48.4,48.4)\kV$ 为附件中位电压。于是
\[
\sum_iw_{i,0}\left(\frac{U_i}{U_{i,0}}\right)^2
=\frac{\sum_iU_i^2}{\sum_iU_{i,0}^2},
\]
即权重只是把 D--A 电压平方代理规范到参考点，不额外引入人为场序系数。计算中令
\[
\bm w\sim\operatorname{Dirichlet}(120\bm w_0),
\]
以传播等效假设与运行波动带来的不确定性；因此它是透明的代理权重，而非分场实测效率。

由删失运行点 $\hat\mu_c$ 和参考入口浓度 $C_{\mathrm{in},0}=37.41\gNm$ 锚定
\begin{equation}
D_0=\log\left(\frac{1000C_{\mathrm{in},0}R_0}{\hat\mu_c}\right)=6.6846,
\label{eq:dref}
\end{equation}
其中 $R_0$ 为参考振打增益。这里 $D_0$ 是附件能提供的最主要排放尺度信息，形状参数仍由文献先验控制。

\subsubsection{温度、积灰与振打修正}

高温会改变烟气密度、离子迁移及粉尘比电阻，实验研究表明温度影响具有明显非线性\cite{xu2015hightemp}。附件温度均值为
\begin{equation}
\Theta_0=\frac{1}{n}\sum_{t=1}^{n}\Theta_t=126.008^\circ\mathrm C,
\label{eq:tempref}
\end{equation}
故将其作为可复算的归一化参考中心，并采用钟形修正
\begin{equation}
g_\Theta(\Theta)=\exp\left[-\kappa_\Theta
\left(\frac{\Theta-\Theta_0}{s_\Theta}\right)^2\right],
\qquad \kappa_\Theta=0.12,\quad s_\Theta=25^\circ\mathrm C.
\label{eq:tempfactor}
\end{equation}
该形式可由 $\log g_\Theta$ 在参考点附近的二阶展开得到，保证修正为正且在参考点等于 1。$\Theta_0$ 来自附件数据；$\kappa_\Theta/s_\Theta^2$ 是文献机理约束下的先验曲率。本文不将 $126.008^\circ\mathrm C$ 解释为设备的普适最优温度，只用于描述本工况区间内偏离参考温区后的有效迁移能力折减。

入口粉尘质量负荷的相对值为
\begin{equation}
L_t=\frac{C_{\mathrm{in},t}Q_t}{C_{\mathrm{in},0}Q_0}.
\label{eq:loadratio}
\end{equation}
若实际振打时刻可得，极板尘负荷可由隐状态方程表示：
\begin{equation}
M_{i,t+1}=(1-\rho_i I_{i,t}^{\mathrm{rap}})M_{i,t}
+\eta_i L_t\bigl(1-\exp(-D_{i,t})\bigr)\Delta t,
\label{eq:duststate}
\end{equation}
其中 $I_{i,t}^{\mathrm{rap}}$ 是振打触发指示量，$\rho_i$ 是清灰比例。附件只给定周期设定值而没有实际触发相位，故无法逐分钟恢复 $M_{i,t}$。对未知相位积分后，以稳态尘负荷指标代替：
\begin{align}
I_M(\bm T,L)&=q_{0.95}(L)\sum_{i=1}^{4}w_i\frac{T_i}{T_{i,0}},\label{eq:dustindex}\\
g_M(\bm T,L)&=\exp\left[-\kappa_M L
\sum_{i=1}^{4}w_i\left(\max\left\{\frac{T_i}{T_{i,0}}-1,0\right\}\right)^2\right].\label{eq:platefactor}
\end{align}
模型用 $I_M\le1.45$ 防止通过无限延长周期虚假节能。

振打使一部分尘饼落入灰斗，同时会造成短时再飞扬\cite{bush1984rapping,neundorfer1981rapping}。其相位积分后的浓度放大因子为
\begin{equation}
R_{\mathrm{rap}}=1+aL\sum_{i=1}^{4}w_i
\left(\frac{T_i}{T_{i,0}}\right)^{1.20}
+b\sum_{i=1}^{4}w_i
\left(\frac{T_i}{T_{i,0}}\right)^{-0.50}.
\label{eq:rapfactor}
\end{equation}
第一项表示周期越长，单次积灰量和再飞扬幅度越大；第二项表示周期过短会增加脉冲出现频率。最终情景浓度为
\begin{equation}
C_{\mathrm{out}}=1000C_{\mathrm{in}}\exp(-D)R_{\mathrm{rap}}\exp(\varepsilon),
\qquad \varepsilon\sim\mathcal N(0,\sigma_\varepsilon^2).
\label{eq:fulltwin}
\end{equation}

\subsubsection{不确定性传播}

对每个工况，从该工况的分钟样本有放回抽取 $(\Theta,C_{\mathrm{in}},Q)$；$D_0$ 采用截断正态扰动，$\bm w$ 采用 Dirichlet 扰动，其余正参数采用对数正态扰动。机会约束的基准情景采用校准后的后验宽度：$D_0$ 相对标准差 2\%，形状参数相对标准差 10\%，对数残差标准差 0.01；另将关键机理系数单独扰动 $\pm30\%$ 作认知不确定性压力测试。于是同一控制策略 $\bm x=(U_1,\ldots,U_4,T_1,\ldots,T_4)$ 对应浓度分布而非单一点预测。

\subsection{时间正则化典型工况划分}

对原始分钟数据保留完整时间顺序。题目二突出入口浓度和温度的波动，故将二者作为工况主轴；考虑到流量同时决定停留时间和入口质量负荷，再把流量作为辅助描述量。首先以 15 分钟居中中位数平滑三个变量，构造
\begin{equation}
\bm z_t=\left(\Theta_t,\log C_{\mathrm{in},t},\log Q_t\right)^\top.
\end{equation}
再用各维中位数与四分位距进行稳健标准化，以降低高温异常对尺度的影响。候选状态数 $K=3,\ldots,8$，每个候选拟合全协方差高斯混合模型
\begin{equation}
p(\bm z_t)=\sum_{k=1}^{K}\pi_k\mathcal N(\bm z_t\mid\bm\mu_k,\bm\Sigma_k).
\label{eq:gmm}
\end{equation}

为抑制分钟级标签跳变，定义发射代价 $d_{tk}=-\log p(k\mid\bm z_t)$，用动态规划求解
\begin{equation}
\min_{s_1,\ldots,s_n}\left[\sum_{t=1}^{n}d_{t,s_t}
+\lambda\sum_{t=2}^{n}\mathbb I(s_t\ne s_{t-1})\right],
\qquad \lambda=8.
\label{eq:viterbi}
\end{equation}
驻留不足 15 分钟的短段合并至相邻低代价状态。有效模型需满足每个状态占比至少 1\%、中位驻留至少 15 分钟；在此基础上选择 BIC 最小者，若 BIC 差不超过 10，则取状态数较小者。

\subsection{能耗代理模型}

电场功率在工作区间内与电压平方近似相关；振打机构的平均功率与动作频率 $1/T_i$ 相关，多电场节能控制研究也支持按场分配能量的思路\cite{li2013ga}。附件字段 $P_{\mathrm{total}}$ 的单位为 kW，故本文先建立总功率模型；当比较时段相同时，最小化功率与最小化电量 $E=P\Delta t$ 等价。建立
\begin{equation}
P(\bm x)=\beta_0+\sum_{i=1}^{4}\beta_iU_i^2
+\sum_{i=1}^{4}\gamma_i\frac{1}{T_i}+\epsilon_P.
\label{eq:power}
\end{equation}
采用稳健回归减弱偶发功率扰动影响，并以逐日留一交叉验证代替随机分钟切分，避免相邻时间点同时进入训练与测试造成泄漏。模型只在历史操作包络及问题 4 的 3\% 局部扩展中使用。

\subsection{95\%机会约束协同优化}

对工况 $r$ 和限值 $L\in\{10,5\}$，优化问题为
\begin{align}
\min_{\bm x_r}\quad & P(\bm x_r),\label{eq:obj}\\
\text{s.t.}\quad & \Prob\left(C_{\mathrm{out}}(\bm x_r,\bm\xi_r)\le L\right)\ge0.95,\label{eq:chance}\\
& I_M(\bm x_r)\le1.45,\label{eq:dustconstraint}\\
& \bm x_{\min}\le\bm x_r\le\bm x_{\max},\label{eq:bounds}
\end{align}
其中 $\bm\xi_r$ 包含工况波动、机理参数和残差不确定性。训练时采用 97.5\% 设计分位预留有限样本裕量；用差分进化执行全局搜索，再用 SLSQP 局部精修。罚函数写为
\begin{equation}
J(\bm x)=P(\bm x)+\Lambda\left[
\max\left(\frac{q_{0.975}(C)}{L}-1,0\right)^2
+\max\left(\frac{I_M}{1.45}-1,0\right)^2\right].
\label{eq:penalty}
\end{equation}
最终策略用另一随机种子生成的 10,000 个独立情景检查 $q_{0.95}(C)\le L$。若首次独立检查未通过，则只沿历史电压上界方向寻找最小安全补偿，并在第二组情景上以 97.5\% 设计分位精修；报告值再由第三组独立 10,000 情景审计。该补偿不改变工况或设备边界，且在结果表中单独标记。算法固定基础种子为 2026，以确保可复现。

\subsection{调节优先级指标}

在最优策略附近，对每个可调量作 1\% 单变量反事实扰动，定义单位新增电耗的 95\% 分位减排收益
\begin{equation}
E_j=\frac{q_{0.95}(C\mid\bm x^*)-q_{0.95}(C\mid\bm x^*+\Delta x_j\bm e_j)}
{P(\bm x^*+\Delta x_j\bm e_j)-P(\bm x^*)}.
\label{eq:priority}
\end{equation}
若变量已在上界，则使用反向差分并统一换算为正向收益。$E_j$ 越大，说明单位电耗带来的风险分位下降越多；但若尘负荷约束活跃，振打周期具备更高的安全优先级，即使其直接 $E_j$ 较小。
